\section{Introduction}
\label{sec:introduction}

A multi-horizon forecaster issues predictions before their targets arrive. Once a block is complete, its errors can guide the next prediction. With a fixed base forecaster, an auxiliary correction must decide which past errors to retain between forecasts and how much state to devote to them.

Residual feedback has been used to correct traffic forecasts~\cite{rescal2022}, while TEFL~\cite{tefl2026} feeds a fully observed multi-step residual vector to a learned adapter trained jointly with the base model. Other fixed-forecaster approaches learn an auxiliary online forecast from Fourier features of observation history~\cite{elf2025}, correct a forecast using recent target statistics~\cite{cosa2026}, or calibrate input and output spectra~\cite{fac2026}. These methods use different information about the recent stream. We study a compressed representation of the base forecaster's own errors that keeps the most recent error explicitly available.

Keeping a complete residual block preserves every horizon-specific error but requires state proportional to its length. A few transform coefficients require less state but lose individual values. Two blocks can therefore have identical retained discrete cosine transform (DCT) coefficients and different final errors. That endpoint is the latest observed error when the next nonoverlapping block is forecast. We test whether retaining it alongside the compressed summary and current base forecast improves correction.

We propose Endpoint-Preserving Online Correction (EPOC) for a fixed multi-horizon forecaster. EPOC retains low-order DCT coefficients of the preceding completed base-residual block and its uncompressed final value. For each channel, it supplies that endpoint as a shared feature to separate online ridge regressions for the retained correction components. Each regression also receives the corresponding coefficient of the current base forecast. The fitted components reconstruct a correction in the same fixed DCT subspace, and a blending weight derived from completed blocks combines it with the base forecast. Updates occur after the full target block has arrived; the retained auxiliary arrays include regression statistics, the residual summary, the fixed DCT basis, and blending state.

We evaluate three questions. First, how do correction accuracy and retained auxiliary arrays compare with other online methods? Second, what do the endpoint and the other input components contribute at a fixed output dimension? Third, how does the number of retained components affect the accuracy--storage choice? The main comparison applies the uncorrected base, five correction comparators, and EPOC to eight multivariate series, DLinear and PatchTST bases, three seeds, and two base-training variants, yielding 96 matched fixed-base conditions. Equal-size first, middle, and mean residual summaries isolate the position and type of retained scalar; input-removal and projected-endpoint controls examine the complete regression input. A separate experiment compares the correction with TEFL-style residual adapters on the same jointly trained bases.

Across the 96 fixed-base conditions, EPOC reduces mean squared error (MSE) by an average of 15.40\% and mean absolute error (MAE) by 9.35\% relative to the uncorrected base, while retaining a median of 6,352 B of auxiliary numerical arrays. The endpoint yields 1.65--2.20\% lower paired MSE than the three equal-size residual summaries. Full ELF achieves a larger mean MSE reduction from the uncorrected base, 19.29\%, with 474,048 median retained bytes. EPOC retains about one seventy-fifth of ELF's auxiliary state, with a 3.89-percentage-point smaller mean MSE reduction. On jointly trained bases, it lowers MSE by 16.69--20.15\% relative to the globally blended TEFL-style adapter applied to the same base.

The contributions are as follows:
\begin{enumerate}
\item We define a compact retained state that pairs truncated residual DCT coefficients with an uncompressed endpoint and uses current-forecast coefficients in channel- and component-wise online correction.
\item We evaluate the retained-information choice through equal-size residual-summary controls, a coefficient-reconstructed endpoint, and input ablations under matched base forecasts and correction dimensions.
\item We compare accuracy and retained arrays over 96 fixed-base conditions, jointly trained adapters, and component and update settings.
\end{enumerate}

Section~\ref{sec:related_work} reviews related correction and adaptation methods. Section~\ref{sec:method} defines the forecast--feedback order and retained state; Section~\ref{sec:experimental_setup} gives the evaluation protocol. Section~\ref{sec:results} presents the results, followed by discussion and conclusions in Sections~\ref{sec:discussion} and~\ref{sec:conclusion}.
